\documentclass{article}
\usepackage{hyperref}
\usepackage{Style}

\nocite{*} % Comentar si quiero citar
%\addbibresource{bibliografia.bib} % Quitar el comentado si quiero usar bibliografia

\begin{document}

\begin{minipage}{2.5cm}
    \includegraphics[width=2cm]{imagen_puc.jpg}
\end{minipage}
\begin{minipage}{12cm}
    {\sc Pontificia Universidad Católica de Chile\\
    Facultad de Matemáticas\\
    Departamento de Matemática\\
    Profesor: Manuel Cabezas -- Ayudante: Benjamín Mateluna}
\end{minipage}
\vspace{2ex}

{\centerline{\bf Introducción a la Geometría - MAT1304}
\centerline{\bf Ayudantía 13}}
\centerline{\bf 23 de abril de 2026}

\vspace{1cm}
\noindent\textbf{Problema 1.} Demuestre la identidad
\begin{equation*}
    tan(2\alpha)+sec(2\alpha)=\frac{cos(\alpha)+sen(\alpha)}{cos(\alpha)-sen(\alpha)}
\end{equation*}

\vspace{5mm}
\noindent\textbf{Problema 2.} Pruebe que $cos(3\alpha)=4cos^{3}(\alpha)-3cos(\alpha)$ para todo
$\alpha\in\R$.

\vspace{5mm}
\noindent\textbf{Problema 3.} Sea $\triangle ABC$ un triángulo con $\angle ACB=\frac{\pi}{4}$ y
$\angle ABC=\frac{\pi}{3}$. Sea $D\in\overline{BC}$ tal que $3\overline{BD}=\overline{CD}$. 
Encuentre el valor de
\begin{equation*}
    \frac{sen(\angle BAD)}{sen(\angle CAD)}
\end{equation*}

\vspace{5mm}
\noindent\textbf{Problema 4.} Sea $\triangle ABC$ y $P$ en el segmento $\overline{BC}$ tal que
$\overline{AP}=p, \overline{BP}=m$ y $\overline{PC}=n$. Demostrar que $a(p^{2}+mn)=b^{2}m+c^{2}n$.

\vspace{1cm}
\noindent\textbf{\underline{Ejercicios Propuestos:}}

\vspace{5mm}
\noindent\textbf{Extra 1.} Encontrar todos los $t\in\left[0,\frac{\pi}{2}\right]$ tales que
$cos(3t)+3cos(t)=\frac{1}{2}$.

\vspace{5mm}
\noindent\textbf{Extra 2.} Considere $\triangle ADC$ rectángulo en $C$. Sea $B\in\overline{AC}$
y $2\angle DAC=\angle DBC$. Si $\overline{BC}=\frac{3}{5}\overline{AB}$. Calcule el área de 
$\triangle ADC$ en función de $\overline{BC}$.

%\printbibliography % Quitar el comentado si quiero usar bibliografia

\end{document}
